\documentclass[11pt,a4paper]{article}
\usepackage[utf8]{inputenc}
\usepackage[T1]{fontenc}
\usepackage[ngerman]{babel}
\usepackage{helvet}
\renewcommand{\familydefault}{\sfdefault}
\usepackage[left=2cm, right=2cm, top=2cm, bottom=2cm]{geometry}
\usepackage{amsmath,amssymb}
\usepackage{array}
\usepackage{booktabs}
\usepackage{tikz}
\usepackage{pgfplots}
\pgfplotsset{compat=1.18}
\usepackage{fancyhdr}
\usepackage{xcolor}
\usepackage{siunitx}
\sisetup{locale = DE}

% Lösungsfarbe
\definecolor{correct}{RGB}{34,139,94}

\pagestyle{fancy}
\fancyhf{}
\fancyhead[L]{\small Physik Klasse 10E}
\fancyhead[R]{\small\bfseries ERWARTUNGSHORIZONT Probeklausur}
\fancyfoot[C]{\small Seite \thepage}
\renewcommand{\headrulewidth}{0.4pt}

\setlength{\parindent}{0pt}
\setlength{\fboxsep}{8pt}

% Aufgaben-Befehl
\newcounter{aufgabe}
\newcommand{\aufgabe}[2]{%
\stepcounter{aufgabe}%
\vspace{0.5cm}%
\noindent\fbox{\parbox{\dimexpr\textwidth-2\fboxsep-2\fboxrule}{%
\textbf{Aufgabe \theaufgabe: #1}\hfill\textbf{#2 BE}%
}}%
\vspace{0.3cm}%
}

\newcommand{\lsg}[1]{\textcolor{correct}{#1}}
\newcommand{\be}[1]{\hfill\textcolor{correct}{\small(#1 BE)}}

\begin{document}

% === KOPFBEREICH ===
\noindent\fbox{\parbox{\dimexpr\textwidth-2\fboxsep-2\fboxrule}{
\centering
{\Large\bfseries Probeklausur: Kinematik}\\[0.2cm]
{\large\bfseries Erwartungshorizont}\\[0.2cm]
{\small Klasse 10E \quad|\quad 90 Minuten \quad|\quad 48 BE}
}}

\vspace{0.4cm}

% === NOTENSPIEGEL (15-NP Sek II MV) ===
\begin{center}
\small
\begin{tabular}{|c|c|c|c|c|c|c|c|c|c|c|c|c|c|c|c|}
\hline
\textbf{NP} & 15 & 14 & 13 & 12 & 11 & 10 & 9 & 8 & 7 & 6 & 5 & 4 & 3 & 2 & 1 \\
\hline
\textbf{ab BE} & 46 & 43,5 & 41 & 38,5 & 36 & 34 & 31,5 & 29 & 26,5 & 24 & 22 & 19,5 & 16 & 13 & 10 \\
\hline
\textbf{ab \%} & 95 & 90 & 85 & 80 & 75 & 70 & 65 & 60 & 55 & 50 & 45 & 40 & 33 & 27 & 20 \\
\hline
\end{tabular}
\end{center}

\vspace{0.3cm}

% ═══════════════════════════════════════════════════════════════
% AUFGABE 1: Grundwissen (8 BE, AFB I)
% ═══════════════════════════════════════════════════════════════

\aufgabe{Grundwissen Kinematik}{8}

\begin{enumerate}
\item[a)] \lsg{Die Fläche unter einem $v$-$t$-Graphen stellt den \textbf{zurückgelegten Weg $s$} dar.} \be{1}

\item[b)] \lsg{Die Steigung im $v$-$t$-Diagramm stellt die \textbf{Beschleunigung $a$} dar.} \be{1}

\item[c)] Umrechnung: \be{2}

\lsg{$72\,\mathrm{km/h} = 72 \div 3{,}6 = 20\,\mathrm{m/s}$}

\lsg{$36\,\mathrm{km/h} = 36 \div 3{,}6 = 10\,\mathrm{m/s}$}

{\small\color{correct} Je 1 BE pro richtige Umrechnung}

\item[d)] Tabelle ergänzen: \be{4}

\begin{center}
\begin{tabular}{|p{6cm}|p{3cm}|p{5cm}|}
\hline
\textbf{Bewegungsart} & \textbf{Diagramm} & \textbf{Graphenform} \\
\hline
Gleichförmige Bewegung & $s$-$t$ & \lsg{Ansteigende Gerade} \\
\hline
Gleichmäßig beschleunigte Bewegung & $s$-$t$ & \lsg{Nach oben geöffnete Parabel} \\
\hline
Abbremsen (negative Beschleunigung) & $v$-$t$ & \lsg{Fallende Gerade} \\
\hline
Stillstand & $v$-$t$ & \lsg{Waagerechte Linie bei $v=0$} \\
\hline
\end{tabular}
\end{center}

{\small\color{correct} Je 1 BE pro richtige Zuordnung}
\end{enumerate}

% ═══════════════════════════════════════════════════════════════
% AUFGABE 2: Brunnenexperiment (6 BE)
% ═══════════════════════════════════════════════════════════════

\aufgabe{Freier Fall -- Brunnenexperiment}{6}

\begin{enumerate}
\item[a)] Umstellen nach $t$: \be{2}

\lsg{Ausgangsformel: $h = \dfrac{1}{2}g \cdot t^2$}\\[0.2cm]
\lsg{Umstellen nach $t^2$: $t^2 = \dfrac{2h}{g}$} \hfill(1 BE)\\[0.2cm]
\lsg{Wurzel ziehen: $t = \sqrt{\dfrac{2h}{g}}$} \hfill(1 BE)

\item[b)] Berechnung der Brunnentiefe: \be{2}

\lsg{$h = \dfrac{1}{2} \cdot g \cdot t^2 = \dfrac{1}{2} \cdot 10\,\mathrm{m\,s^{-2}} \cdot (3\,\mathrm{s})^2$} \hfill(1 BE)\\[0.2cm]
\lsg{$h = \dfrac{1}{2} \cdot 10 \cdot 9\,\mathrm{m} = 45\,\mathrm{m}$} \hfill(1 BE)

\item[c)] Erklärung Schallausbreitung: \be{2}

\lsg{Die gemessene Zeit $t_{\text{gemessen}}$ enthält sowohl die Fallzeit als auch die Zeit für den Schall, nach oben zu gelangen.} \hfill(1 BE)\\[0.2cm]
\lsg{Da die tatsächliche Fallzeit kürzer ist als gemessen, ist die tatsächliche Brunnentiefe \textbf{kleiner} als berechnet.} \hfill(1 BE)
\end{enumerate}

\newpage

% ═══════════════════════════════════════════════════════════════
% AUFGABE 3: Bremsvorgang Fahrrad (8 BE)
% ═══════════════════════════════════════════════════════════════

\aufgabe{Bremsvorgang eines Fahrrads}{8}

\begin{enumerate}
\item[a)] \lsg{$v_0 = 36\,\mathrm{km\,h^{-1}} = 36 \div 3{,}6 = 10\,\mathrm{m\,s^{-1}}$} \be{1}

\item[b)] Bremszeit: \be{2}

\lsg{Formel: $t = \dfrac{v_0}{|a|}$} \hfill(1 BE)\\[0.2cm]
\lsg{Einsetzen: $t = \dfrac{10\,\mathrm{m\,s^{-1}}}{5\,\mathrm{m\,s^{-2}}} = 2\,\mathrm{s}$} \hfill(1 BE)

\item[c)] Bremsweg: \be{2}

\lsg{Formel: $s = \dfrac{v_0^2}{2|a|}$} \hfill(1 BE)\\[0.2cm]
\lsg{Einsetzen: $s = \dfrac{(10\,\mathrm{m\,s^{-1}})^2}{2 \cdot 5\,\mathrm{m\,s^{-2}}} = \dfrac{100\,\mathrm{m^2\,s^{-2}}}{10\,\mathrm{m\,s^{-2}}} = 10\,\mathrm{m}$} \hfill(1 BE)

\item[d)] Beurteilung: \be{3}

\lsg{\textbf{Ja}, der Fahrradfahrer kann rechtzeitig stoppen.} \hfill(1 BE)\\[0.2cm]
\lsg{Begründung: Der Bremsweg beträgt 10 m, der Hund ist 12 m entfernt.} \hfill(1 BE)\\[0.2cm]
\lsg{$12\,\mathrm{m} - 10\,\mathrm{m} = 2\,\mathrm{m}$ Sicherheitsabstand verbleiben.} \hfill(1 BE)
\end{enumerate}

% ═══════════════════════════════════════════════════════════════
% AUFGABE 4: Straßenbahnfahrt (12 BE)
% ═══════════════════════════════════════════════════════════════

\aufgabe{Straßenbahnfahrt}{12}

\begin{enumerate}
\item[a)] $v$-$t$-Diagramm zeichnen: \be{3}

\begin{center}
\begin{tikzpicture}
\begin{axis}[
    width=13cm, height=5cm,
    xlabel={$t$ in s},
    ylabel={$v$ in m/s},
    xmin=0, xmax=18,
    ymin=0, ymax=15,
    grid=both,
    grid style={gray!30},
    major grid style={gray!50},
    axis lines=left,
    xtick={0,4,8,12,16},
    ytick={0,4,8,12},
    every axis label/.style={font=\small},
]
% Phase 1: Anfahren (0-4s), 0→12 m/s
\addplot[very thick, correct] coordinates {(0,0) (4,12)};
% Phase 2: Gleichförmig (4-12s), 12 m/s
\addplot[very thick, correct] coordinates {(4,12) (12,12)};
% Phase 3: Bremsen (12-16s), 12→0 m/s
\addplot[very thick, correct] coordinates {(12,12) (16,0)};
% Markierungspunkte
\addplot[only marks, mark=*, correct, mark size=2pt] coordinates {(0,0) (4,12) (12,12) (16,0)};
\end{axis}
\end{tikzpicture}
\end{center}

\textbf{Bewertungskriterien:}\\
\lsg{Achsenbeschriftung korrekt ($t$ in s, $v$ in m/s)} \hfill(1 BE)\\
\lsg{Korrekte Werte: Anfahren 0--4\,s auf 12\,m/s, Fahrt 4--12\,s bei 12\,m/s, Bremsen 12--16\,s auf 0} \hfill(1 BE)\\
\lsg{Korrekte Graphenform: Ansteigende Gerade, Waagerechte, Fallende Gerade} \hfill(1 BE)

\item[b)] Berechnung des Gesamtweges (Flächenansatz): \be{4}

\lsg{\textbf{Phase 1 (Dreieck):} $s_1 = \dfrac{1}{2} \cdot 4\,\mathrm{s} \cdot 12\,\mathrm{m\,s^{-1}} = 24\,\mathrm{m}$} \hfill(1 BE)\\[0.2cm]
\lsg{\textbf{Phase 2 (Rechteck):} $s_2 = 8\,\mathrm{s} \cdot 12\,\mathrm{m\,s^{-1}} = 96\,\mathrm{m}$} \hfill(1 BE)\\[0.2cm]
\lsg{\textbf{Phase 3 (Dreieck):} $s_3 = \dfrac{1}{2} \cdot 4\,\mathrm{s} \cdot 12\,\mathrm{m\,s^{-1}} = 24\,\mathrm{m}$} \hfill(1 BE)\\[0.2cm]
\lsg{\textbf{Gesamtweg:} $s_{\text{ges}} = s_1 + s_2 + s_3 = 24\,\mathrm{m} + 96\,\mathrm{m} + 24\,\mathrm{m} = 144\,\mathrm{m}$} \hfill(1 BE)

{\small\color{correct} Alternativ: Gesamtfläche als Trapez: $s = \frac{1}{2}(4 + 12) \cdot 12 + \frac{1}{2} \cdot 4 \cdot 12 = 96 + 48 = 144\,\mathrm{m}$}

\item[c)] $s$-$t$-Diagramm zeichnen: \be{5}

\begin{center}
\begin{tikzpicture}
\begin{axis}[
    width=13cm, height=6cm,
    xlabel={$t$ in s},
    ylabel={$s$ in m},
    xmin=0, xmax=18,
    ymin=0, ymax=160,
    grid=both,
    grid style={gray!30},
    major grid style={gray!50},
    axis lines=left,
    xtick={0,4,8,12,16},
    ytick={0,24,48,72,96,120,144},
    every axis label/.style={font=\small},
]
% Phase 1: Parabel (beschleunigt) - s = 1.5*t^2
\addplot[very thick, correct, samples=50, domain=0:4] {1.5*x^2};
% Phase 2: Gerade (gleichförmig) - von (4,24) bis (12,120)
\addplot[very thick, correct] coordinates {(4,24) (12,120)};
% Phase 3: Parabel (bremst) - s = 144 - 1.5*(16-t)^2
\addplot[very thick, correct, samples=50, domain=12:16] {144 - 1.5*(16-x)^2};
% Markierungspunkte
\addplot[only marks, mark=*, correct, mark size=2pt] coordinates {(0,0) (4,24) (12,120) (16,144)};
% Gestrichelte vertikale Phasengrenzen
\draw[gray!70, dashed, thick] (axis cs:4,0) -- (axis cs:4,160);
\draw[gray!70, dashed, thick] (axis cs:12,0) -- (axis cs:12,160);
\end{axis}
\end{tikzpicture}
\end{center}

\textbf{Bewertungskriterien:}\\
\lsg{Achsenbeschriftung korrekt ($t$ in s, $s$ in m mit sinnvollen Werten)} \hfill(1 BE)\\
\lsg{Phase 1: Nach oben gekrümmte Parabel (Beschleunigung), endet bei $s = 24\,\mathrm{m}$} \hfill(1,5 BE)\\
\lsg{Phase 2: Ansteigende Gerade (gleichförmig), endet bei $s = 120\,\mathrm{m}$} \hfill(1 BE)\\
\lsg{Phase 3: Nach unten gekrümmte Parabel (Bremsen), endet bei $s = 144\,\mathrm{m}$} \hfill(1,5 BE)

\end{enumerate}

\newpage

% ═══════════════════════════════════════════════════════════════
% AUFGABE 5: Radfahrer überholt Jogger (8 BE)
% ═══════════════════════════════════════════════════════════════

\aufgabe{Radfahrer überholt Jogger}{8}

\begin{enumerate}
\item[a)] Umrechnung: \be{1}

\lsg{$v_{\text{Jogger}} = 9\,\mathrm{km\,h^{-1}} = 9 \div 3{,}6 = 2{,}5\,\mathrm{m\,s^{-1}}$}\\
\lsg{$v_{\text{Radfahrer}} = 18\,\mathrm{km\,h^{-1}} = 18 \div 3{,}6 = 5\,\mathrm{m\,s^{-1}}$}

\item[b)] Relativgeschwindigkeit: \be{2}

\lsg{Formel: $v_{\text{rel}} = v_{\text{Radfahrer}} - v_{\text{Jogger}}$} \hfill(1 BE)\\[0.2cm]
\lsg{Einsetzen: $v_{\text{rel}} = 5\,\mathrm{m\,s^{-1}} - 2{,}5\,\mathrm{m\,s^{-1}} = 2{,}5\,\mathrm{m\,s^{-1}}$} \hfill(1 BE)

\item[c)] Zeit bis zum Erreichen: \be{2}

\lsg{Formel: $t = \dfrac{\Delta s}{v_{\text{rel}}}$} \hfill(1 BE)\\[0.2cm]
\lsg{Einsetzen: $t = \dfrac{50\,\mathrm{m}}{2{,}5\,\mathrm{m\,s^{-1}}} = 20\,\mathrm{s}$} \hfill(1 BE)

\item[d)] Analyse des Überholvorgangs: \be{3}

\lsg{In $t = 20\,\mathrm{s}$ legt der Radfahrer zurück:}\\
\lsg{$s_{\text{Rad}} = v_{\text{Rad}} \cdot t = 5\,\mathrm{m\,s^{-1}} \cdot 20\,\mathrm{s} = 100\,\mathrm{m}$}\\[0.2cm]
\lsg{Der Radfahrer hat den Jogger nach 20\,s genau an der Kurve \textbf{erreicht}.} \hfill(1 BE)\\[0.2cm]
\lsg{Für sichere Überholung (+ 10\,m Sicherheitsabstand) braucht er zusätzlich:}\\
\lsg{$t_{\text{extra}} = \dfrac{10\,\mathrm{m}}{2{,}5\,\mathrm{m\,s^{-1}}} = 4\,\mathrm{s}$}\\
\lsg{Gesamt: $t_{\text{ges}} = 20\,\mathrm{s} + 4\,\mathrm{s} = 24\,\mathrm{s}$}\\[0.2cm]
\lsg{Weg des Radfahrers in 24\,s: $s = 5\,\mathrm{m\,s^{-1}} \cdot 24\,\mathrm{s} = 120\,\mathrm{m} > 100\,\mathrm{m}$} \hfill(1 BE)\\[0.2cm]
\lsg{\textbf{Beurteilung:} Der Radfahrer kann den Jogger \textbf{nicht} vor der Kurve sicher überholen, da er dafür 120\,m bräuchte, aber nur 100\,m zur Verfügung stehen.} \hfill(1 BE)
\end{enumerate}

% ═══════════════════════════════════════════════════════════════
% AUFGABE 6: Kaffeefilter-Experiment (6 BE)
% ═══════════════════════════════════════════════════════════════

\aufgabe{Kaffeefilter-Experiment}{6}

\begin{enumerate}
\item[a)] Beschleunigung nach dem Loslassen: \be{2}

\lsg{Nach dem Loslassen wirkt die Gewichtskraft $F_G = m \cdot g$ auf den Filter.} \hfill(1 BE)\\[0.2cm]
\lsg{Da der Luftwiderstand anfangs gering ist (weil $v$ noch klein), überwiegt die Gewichtskraft und der Filter wird beschleunigt.} \hfill(1 BE)

\item[b)] Konstante Fallgeschwindigkeit: \be{2}

\lsg{Mit zunehmender Geschwindigkeit wächst der Luftwiderstand $F_W$ (er ist proportional zu $v^2$).} \hfill(1 BE)\\[0.2cm]
\lsg{Wenn $F_W = F_G$ gilt (Kräftegleichgewicht), ist die resultierende Kraft null, also $a = 0$. Der Filter fällt dann mit konstanter Grenzgeschwindigkeit.} \hfill(1 BE)

\item[c)] Gefalteter Filter fällt schneller: \be{2}

\lsg{Der gefaltete Kaffeefilter hat eine kleinere Querschnittsfläche $A$ als der ungefaltete.} \hfill(1 BE)\\[0.2cm]
\lsg{Da $F_W \sim A$, ist der Luftwiderstand bei gleicher Geschwindigkeit geringer. Das Kräftegleichgewicht $F_W = F_G$ wird erst bei höherer Geschwindigkeit erreicht -- die Grenzgeschwindigkeit ist größer.} \hfill(1 BE)
\end{enumerate}

\vspace{0.5cm}
\hrule
\vspace{0.5cm}

% === AFB-ZUSAMMENFASSUNG ===

\begin{center}
\fbox{\parbox{0.9\textwidth}{
\centering
\textbf{Zusammenfassung der Anforderungsbereiche}\\[0.3cm]
\begin{tabular}{|l|c|c|c|}
\hline
\textbf{AFB} & \textbf{Teilaufgaben} & \textbf{BE} & \textbf{Anteil} \\
\hline
I (Reproduktion) & 1a-d, 2a, 3a, 5a & 12 & 25\% \\
\hline
II (Anwendung) & 2b, 3b-c, 4a-c, 5b-c, 6a & 20 & 42\% \\
\hline
III (Transfer) & 2c, 3d, 4d, 5d, 6b-c & 16 & 33\% \\
\hline
\textbf{Gesamt} & & \textbf{48} & 100\% \\
\hline
\end{tabular}
}}
\end{center}

\vspace{0.3cm}

\begin{center}
\textit{--- Ende des Erwartungshorizonts ---}
\end{center}

\end{document}
